\section[Integration and Fourier transform on $\mR^m_q$]{Integration and Fourier transform on $\boldsymbol{\mR^m_q}$}

\subsection[Integration over the quantum sphere and induced integration on $\mR^m_q$]{Integration over the quantum sphere and induced integration on $\boldsymbol{\mR^m_q}$}
\label{DefFourier}

First we show how the Howe dual pair $(O_q(m),\cU_q(\mathfrak{sl}_2))$ uniquely characterizes the integration over the quantum sphere from~\cite{MR1408120}.

\begin{theorem}
\label{Pizzth}
The unique $($up to a multiplicative constant$)$ linear functional on $\cP$ invariant under the co-action of $O_q(m)$ and satisfying $\int_{\mS^{m-1}_q}\bold{x}^2R=\int_{\mS^{m-1}_q}R$ is given by the Pizzetti formula
\begin{gather*}
%\label{Pizzetti}
\int_{\mS^{m-1}_q}R= \sum_{k=0}^\infty\frac{2\left(\Gamma_{q^2}(\frac{1}{2})\right)^{m}}{\mu^{2k}[k]_{q^2}!
\Gamma_{q^2}(k+\frac{m}{2})}\big(\Delta^kR\big)(0)\qquad\mbox{for}\quad R\in\cP.
\end{gather*}
\end{theorem}

\begin{proof}
The Fischer decomposition in Lemma~\ref{Fischer} implies that the integration on~$\cP$ is uniquely determined if it is determined on each of the blocks $\bold{x}^{2l}\cS_k$. Since these blocks are irreducible $\cU_q(\mathfrak{so}(m))$-representations (or irreducible $O_q(m)$-corepresentations) the integration should be zero on each such block which is not one dimensional. This implies that the integration can only have non-zero values on the elements~$\bold{x}^{2l}$. The second property then implies $\int_{\mS^{m-1}_q}\bold{x}^{2l}=\int_{\mS^{m-1}_q}1$, which shows that the integration is uniquely determined up to the constant~$\int_{\mS^{m-1}_q}1$. It is easily checked that the Pizzetti formula satisf\/ies the conditions.
\end{proof}

We chose the normalization such that $\int_{\mS^{m-1}_q}1
=\frac{2\left(\Gamma_{q^2}(\frac{1}{2})\right)^{m}}{\Gamma_{q^2}(\frac{m}{2})}$.
This quantum sphere integration can be expressed symbolically using the f\/irst $q$-Bessel-function~\eqref{Bessel1},
\begin{gather}
\label{PizzBessel}
\int_{\mS^{m-1}_q}R = 2\left(\Gamma_{q^2}\left(\frac{1}{2}\right)\right)^{m}\left(\frac{J^{(1)}_{\frac{m}{2}-1}
\left(\frac{1+q}{\mu}\sqrt{-\Delta}|q^2\right)}{\left(\sqrt{-\Delta}/\mu\right)^{\frac{m}{2}-1}}R\right)(0).
\end{gather}

The following lemma will be important for the sequel.
\begin{lemma}\label{xS}
The Fischer decomposition $($see Lemma~{\rm \ref{Fischer})} of $x^jS_k(\bold{x})$ with $S_k\in\cS_k$ is given by
\begin{gather*}
x^jS_k=\left(x^jS_k-\bold{x}^2\frac{1}{\mu[k+m/2-1]_{q^2}}\partial^jS_k\right)+\bold{x}^2\left(\frac{1}{\mu[k+m/2-1]_{q^2}}\partial^jS_k\right)
\end{gather*}
with
\begin{gather*}
\left(x^jS_k-\bold{x}^2\frac{1}{\mu[k+m/2-1]_{q^2}}\partial^jS_k\right)\in\cS_{k+1} \!\!\!\qquad \mbox{and} \qquad\!\!\! \left(\frac{1}{\mu[k+m/2-1]_{q^2}}\partial^jS_k\right)\in\cS_{k-1}.
\end{gather*}
\end{lemma}
\begin{proof}\sloppy
Since $\Delta$ and $\partial^j$ commute $\left(\frac{1}{\mu[k+m/2-1]_{q^2}}\partial^jS_k\right)\in\cS_{k-1}$. Equations \eqref{Laplx}, \eqref{commLaplxkwad} and~\eqref{Eharm} yield
\begin{gather*}
\Delta\left(x^jS_k-\bold{x}^2\frac{1}{\mu[k+m/2-1]_{q^2}}\partial^jS_k\right)=
\mu\partial^jS_k-\mu^2E\frac{1}{\mu[k+m/2-1]_{q^2}}\partial^jS_k=0.
\end{gather*}
This can also be calculated by using the projection operator from $\cP_k$ onto $\cS_k$ developed \linebreak in~\cite{MR2006509}.
\end{proof}


Before we use the quantum sphere integration to construct integration on $\mR^m_q$ along the lines of~\cite{MR1303081} and~\cite{MR1239953}, the following lemma should be considered.
\begin{lemma}
\label{intexpinf}
For the $q$-integration in equation \eqref{gammaint} with $\gamma\in\mR^+$, the expression
\begin{gather*}
\int_0^{\gamma\cdot\infty}E_q(-t)d_qt
\end{gather*}
is infinite unless $\gamma=q^{j}\frac{1}{1-q}$ for some $j\in\mZ$. In that case it reduces to $\int_0^{\frac{1}{1-q}}E_q(-t)d_qt$.
\end{lemma}

\begin{proof}
Property \eqref{aflexp} can be rewritten as $E_q(t)=E_q(qt)\left[1+(1-q)t\right]$, which implies
\begin{gather*}
E_q\big({-}q^{-k}\gamma\big) = E_q\big({-}q^{1-k}\gamma\big)\big[1-(1-q)q^{-k}\gamma\big].
\end{gather*}

First we assume $E_q(-q^{1-k}\gamma)$ is never zero for $k\in\mZ$. The equation above then yields $\left|E_q(-q^{-k}\gamma)\right|>\left|E_q(-q^{1-k}\gamma)\right|q$ and ${\rm Sign}(E_q(-q^{-k}\gamma))\not={\rm Sign}(E_q(-q^{1-k}\gamma))$ for $k\ge N\in\mN$ with $N>\frac{\ln\left(\frac{1+q}{\gamma(1-q)}\right)}{\ln(1/q)}$. This implies that $\sum\limits_{k=N}^\infty E_q(-q^{-k}\gamma)q^{-k}$ will not converge, therefore the integral~\eqref{gammaint} will not exist.

If $E_q(-q^{1-k}\gamma)=0$ holds for some $k\in\mZ$, the unicity of the zeroes in equation~\eqref{zeroes} implies that $\gamma=\frac{q^{k-1-l}}{1-q}$ for some $k\in\mZ$ and $l\in\mN$, or $\gamma=\frac{q^{j}}{1-q}$ for some $j\in\mZ$. The integration then reduces to the proposed expression for the same reason as in equation~\eqref{intinfeind}.
\end{proof}


Now we can construct integration on $\mR^m_q$ from the quantum sphere integration according to the principle of Gaussian-induced integration. The two dif\/ferential calculi lead to dif\/ferent Gaussians and integrations. For constants $\alpha,\beta\in\mR^+$, the relations
\begin{gather*}
\partial^je_{q^2}\big({-}\alpha\bold{x}^2\big)=-\alpha\mu x^je_{q^2}\big({-}\alpha\bold{x}^2\big)\qquad \mbox{and}\qquad \overline{\partial^j} E_{q^2}\big({-}\beta\bold{x}^2\big)=-\beta q^{m-2}\mu x^jE_{q^2}\big({-}\beta\bold{x}^2\big)
\end{gather*}
imply that $e_{q^2}(-\alpha\bold{x}^2)$ and $E_{q^2}(-\beta\bold{x}^2)$ are Gaussians according to equation~\eqref{Gaussian} for the unbarred and barred calculus respectively. Def\/ine the following two integrations for $\gamma,\lambda\in\mR^+$ on polynomials weighted with undetermined radial functions,
\begin{gather*}
\int_{\gamma\cdot\mR^m_q}R_{k}(\bold{x})f\big(\bold{x}^2\big) = \int_0^{\gamma\cdot\infty}d_qr\,r^{m+k-1}f\big(r^2\big)\int_{\mS^{m-1}_q}R_{k}(\bold{x})\qquad \mbox{and}\\
\int_{\mB^m_q(\lambda)}R_{k}(\bold{x})f\big(\bold{x}^2\big) = \int_0^{\lambda}d_qr\,r^{m+k-1}f\big(r^2\big)\int_{\mS^{m-1}_q}R_{k}(\bold{x}).
\end{gather*}
The Taylor expansion of the function $f$ in the origin is assumed to converge on $\mR^+$ for the f\/irst integration and on $[0,\lambda]$ for the second integration. These integrations correspond to the Gaussian-induced integrations.
\begin{theorem}
\label{2indint}
The integrations defined above satisfy
\begin{gather*}
\int_{\gamma\cdot\mR^m_q}\partial^j=0 \qquad \mbox{on}\quad  \cP\otimes e_{q^2}\big({-}\alpha\bold{x}^2\big)\quad \mbox{for}\quad \alpha,\gamma\in\mR^+\qquad\mbox{and}\\
\int_{\mB^m_q\big(\frac{1}{\sqrt{(1-q^2)\beta}}\big)}\overline{\partial^j}=0\qquad \mbox{on}\quad \cP\otimes E_{q^2}\big({-}\beta\bold{x}^2\big)\quad\mbox{for}\quad\beta\in\mR^+.
\end{gather*}
\end{theorem}
\begin{proof}
The f\/irst property is well-known, see e.g.~\cite{MR1239953, MR1408120}. In order to prove the second property we consider the expression
\begin{gather*}
\int_{\mB^m_q(\lambda)}\overline{\partial^j}f\big(\bold{x}^2\big)S_k(\bold{x})
=q^{m-2}\mu\int_{\mB^m_q(\lambda)}\big(\partial_{\bold{x}^2}^{q^{-2}}f\big(\bold{x}^2\big)\big)x^j S_k(\bold{x})+\int_{\mB^m_q(\lambda)}f\big(q^{-2}\bold{x}^2\big)\overline{\partial^j}S_k(\bold{x}).
\end{gather*}
for a spherical harmonic $S_k\in\cS_k$. If $k\not=1$ the right-hand side is always zero, because of Lemma~\ref{xS} and the expression for $\int_{\mS^{m-1}_q}$ in Theorem~\ref{Pizzth}. In case $k=1$ the spherical harmonics are the monomials $x^i$, $i=1,\dots,m$. Therefore we calculate using equation~\eqref{berpx2}, Leibniz rule~\eqref{Leibniz} and equation~\eqref{partint1}
\begin{gather*}
\int_{\mB^m_q(\lambda)}\overline{\partial^j}f\big(\bold{x}^2\big)x^i
=q^{m-2}\mu\int_{\mB^m_q(\lambda)}\big(\partial_{\bold{x}^2}^{q^{-2}}f\big(\bold{x}^2\big)\big)x^j x^i+\int_{\mB^m_q(\lambda)}f\big(q^{-2}\bold{x}^2\big)C^{ji}\\
=\frac{2\left(\Gamma_{q^2}(\frac{1}{2})\right)^{m}C^{ji}}{\Gamma_{q^2}(1+\frac{m}{2})}\left[q^{m-2}\!\int_0^\lambda \!d_qr r^{m+1}\frac{1}{(1+q^{-1})r}\partial^{q^{-1}}_rf\big(r^2\big)+\left[\frac{m}{2}\right]_{q^2}\!\int_0^\lambda \! d_qr r^{m-1}f\big(q^{-2}r^2\big)\right]\\
=\frac{2\left(\Gamma_{q^2}(\frac{1}{2})\right)^{m}C^{ji}}{\Gamma_{q^2}(1+\frac{m}{2})(1+q)}\left[q^{m}\int_0^\lambda d_qr r^{m}\partial^q_rf\big(q^{-2}r^2\big)+\left[m\right]_q\int_0^\lambda d_qr r^{m-1}f\big(q^{-2}r^2\big)\right]\\
=\frac{2\left(\Gamma_{q^2}(\frac{1}{2})\right)^{m}C^{ji}}{\Gamma_{q^2}(1+\frac{m}{2})(1+q)}\left[\lambda^mf\big(q^{-2}\lambda^2\big)-\lim_{r\to0}r^m f\big(q^{-2}r^2\big)\right].
\end{gather*}
When we substitute $f(\bold{x}^2)=\bold{x}^{2l}E_{q^2}(-\beta \bold{x}^2)$ with $l\in\mN$ and $\beta\in\mR^+$ and use equation \eqref{zeroes} we obtain
\begin{gather*}
\int_{\mB^m_q\big(\frac{1}{\sqrt{(1-q^2)\beta}}\big)}\overline{\partial^j}\,\bold{x}^{2l}S_k(\bold{x})\,E_{q^2}\big({-}\beta \bold{x}^2\big) = 0\qquad \forall\, k,l.
\end{gather*}
This proves the second part of the theorem because of the Fischer decomposition in Lem\-ma~\ref{Fischer}.
\end{proof}

\begin{remark}
\label{gengauss}
The properties
\begin{gather*}
\int_{\gamma\cdot\mR^m_q}\overline{\partial^j}=0\qquad \mbox{on}\quad  \cP\otimes e_{q^2}\big({-}\alpha\bold{x}^2\big)\quad \mbox{for}\quad \alpha,\gamma\in\mR^+\qquad\mbox{and}\\
\int_{\mB^m_q\big(\frac{1}{\sqrt{(1-q^2)\beta}}\big)}\partial^j=0\qquad \mbox{on}\quad \cP\otimes E_{q^2}\big({-}\beta\bold{x}^2\big)\quad\mbox{for}\quad\beta\in\mR^+
\end{gather*}
also hold. They do not correspond to Gaussian-induced integration for those calculi in the strict sense. However it is a straightforward generalization, the generalized Gaussian satisf\/ies the relation
\begin{gather*}
-\eta^i_{ j}\partial^jg_{\eta}(\bold{x}) = x^i g_\eta(q\bold{x})
\end{gather*}
(or $-\eta^i_{ j}\overline{\partial^j}g_{\eta}(\bold{x})=x^i g_\eta(q^{-1}\bold{x})$) in stead of equation~\eqref{Gaussian}.
\end{remark}

Since $\int_{\gamma\cdot\mR^m_q}E_{q^2}(-\beta\bold{x}^2)$ will be inf\/inite for general $\gamma$, similar to Lemma~\ref{intexpinf}, the f\/inite integration needs to be used on $\cP\otimes E_{q^2}(-\beta\bold{x}^2)$. As in equation~\eqref{intinfeind} this integration corresponds to a specif\/ic inf\/inite integration,
\begin{gather*}
\label{intbetainf}
\int_{\mB^m_q\big(\frac{1}{\sqrt{(1-q^2)\beta}}\big)} = \int_{\big(\frac{1}{\sqrt{(1-q^2)\beta}}\big)\cdot\mR^m_q}\qquad\mbox{on}\quad \cP\otimes E_{q^2}\big({-}\beta\bold{x}^2\big)\qquad\mbox{for}\quad\beta\in\mR^+.
\end{gather*}
Theorem~\ref{2indint} therefore shows that for integration on $\cP\otimes e_{q^2}(-\alpha\bold{x}^2)$ there is a bigger choice since~$\gamma$ does not depend on $\alpha$. However, it is straightforward to calculate
\begin{gather*}
\int_{\gamma\cdot\mR^m_q}\bold{x}^{2l}S_ke_{q^2}\big({-}\alpha\bold{x}^2\big)=
\delta_{k0}\frac{\Gamma_{q^2}(\frac{m}{2}+l)}{\Gamma_{q^2}(\frac{m}{2})\alpha^lq^{\frac{1}{2}l(l+\frac{m}{2}-1)}}
\int_{\gamma\cdot\mR^m_q}e_{q^2}\big({-}\alpha\bold{x}^2\big),
\end{gather*}
which implies the integrals for dif\/ferent $\gamma$ on $\cP\otimes e_{q^2}(-\alpha\bold{x}^2)$ are proportional to each other.


In the sense of equation~\eqref{Gaussian2} the only dif\/ference between the integrations for dif\/ferent~$\gamma$ is the value $I(g_\eta)$ while the functional~$Z$ does not depend on~$\gamma$. For strict Gaussian-induced integration in the unbarred case, only one choice gives an~$I(g_\eta)$ which is f\/inite.



\subsection[Bochner's relations for the Fourier transform on $\mR^m_q$]{Bochner's relations for the Fourier transform on $\boldsymbol{\mR^m_q}$}\label{boch}

The exponential $\exp_{\hat{R}}(\bold{x}|\bold{y})$ on a quantum space satisf\/ies
\begin{gather}
\label{derFK}
\partial_x^j\exp_{\hat{R}}(\bold{x}|\bold{y}) = \exp_{\hat{R}}(\bold{x}|\bold{y})y^j,
\end{gather}
see \cite{MR1303081, MR1235979, MR2099760}. From now on $\exp_{\hat{R}}(\bold{x}|\bold{y})$ stands for the exponential on $\mR^m_q$. It is uniquely determined from equation~\eqref{derFK} and the normalization $\exp_{\hat{R}}(0|\bold{y})=1$. In order to def\/ine the Fourier transform according to~\cite{MR1303081} the Fourier kernel needs to be combined with the Gaussian induced integrations for the unbarred calculus in Theorem~\ref{2indint} and evaluated on  $\cP\otimes e_{q^2}(-\alpha\bold{x}^2)$. It turns out that the Fourier transform def\/ined by the generalized Gaussian-induced integration for the unbarred calculus in Remark~\ref{gengauss} will lead to interesting properties, see Section~\ref{behave}. This choice also corresponds to the one dimensional theory in equation~\eqref{1dFour}. The Fourier transform can be def\/ined on each space $\cP\otimes E_{q^2}(-\beta \frac{\bold{x}^2}{\mu})$. First we will extend this space. Def\/ine
\begin{gather*}
\overline{\cV}_\beta =\cP\otimes E_{q^2}\left(-\beta \frac{\bold{x}^2}{\mu}\right),\quad\beta\in\mR^+ \qquad \mbox{and} \qquad \cV_\alpha=\cP\otimes e_{q^2}\left(-\alpha\frac{q^2\bold{x}^2}{\mu}\right),\qquad\alpha\in\mR^+.
\end{gather*}
The dif\/ferent spaces $\cV_\alpha$ (or $\overline{\cV}_\beta$) are not necessarily disjunct since equation~\eqref{aflexp} implies
\begin{gather*}
%\label{inclV}
\cV_{q^{-2j}\alpha}\subset\cV_{q^{-2j-2}\alpha} \qquad \mbox{and} \qquad \overline{\cV}_{q^{2j}\beta}\subset\overline{\cV}_{q^{2j+2}\beta}\qquad\mbox{for}\quad j\in\mN.
\end{gather*}
Therefore we def\/ine $\cV_{[\alpha]}=\cup_{j=0}^\infty \cV_{q^{-2j}\alpha}$ and $\overline{\cV}_{[\beta]}=\cup_{j=0}^\infty \overline{\cV}_{q^{2j}\beta}$ for $\alpha,\beta\in\mR^+$. Since, for $j\in\mN$, $\cV_{q^{2j}\alpha}\subset\cV_{\alpha}$,  $\cV_{[\alpha]}$ can also be identif\/ied with $\cup_{j=-\infty}^\infty \cV_{q^{-2j}\alpha}$ and $\overline{\cV}_{[\beta]}$ with $\cup_{j=-\infty}^\infty \overline{\cV}_{q^{2j}\beta}$.

\begin{definition}
\label{defFour1}
The f\/irst Fourier transform $\overline{\cF}^{\pm}_{\mR^m_q}$ on $\mR^m_q$ is a map $\overline{\cV}_{[\beta]}\to\cV_{[1/\beta]}$ for each $\beta\in\mR^+$,
\begin{gather*}
\overline{\cF}^\pm_{\mR^m_q}[f](\bold{y}) = \frac{1+q}{2\mu^{\frac{m}{2}}\left(\Gamma_{q^2}(\frac{1}{2})\right)^{m}}\int_{\frac{1}{\sqrt{(1-q^2)\beta}}\cdot\mR^m_q}f(\bold{x})\exp_{\hat{R}}(\pm i\bold{x}|\bold{y}).
\end{gather*}
\end{definition}
In the subsequent Corollary~\ref{BochFour}, it will be proven that the Fourier transform does indeed map elements of $\overline{\cV}_{[\beta]}$ to $\cV_{[1/\beta]}$. First we need the following technical lemma. For a polynomial~$P$, we def\/ine $[P]^j_l$ by the equation $\sum_l [P]^j_l\partial_x^l=\partial_x^j P-\left[\partial^j_xP\right]$ in ${\rm Dif\/f}(\mR^m_q)$. So in the undeformed Euclidean case, $[P]^j_l=P\delta_l^j$.
\begin{lemma}\label{Svrexp}
For $S_k\in\cS_k$ the relation $\sum_{j,l}[\partial_jS_k(\bold{x})]^j_lx^l=[k]_{q^2}S_k(\bold{x})$ holds with $\partial_j=C_{jk}\partial^k$.
\end{lemma}

\begin{proof}
The lemma is proven by calculating the expression $\sum_j [\partial_x^j[\partial_jS_k(\bold{x})]\bold{x}^2 ]$ in two dif\/ferent ways. The expression is equal to
\begin{gather*}
\sum_{j,l}\left[[\partial_jS_k(\bold{x})]^j_l\partial_x^l\bold{x}^2\right] = \sum_{j,l}[\partial_jS_k(\bold{x})]^j_l\mu x^l.
\end{gather*}
It can also be calculated using formulas \eqref{Laplx} and \eqref{Eharm}
\begin{gather*}
\sum_j\left[\partial_x^j\bold{x}^2[\partial_jS_k(\bold{x})]\right] = \sum_j\mu x^j[\partial_jS_k(\bold{x})]+\sum_jq^2\bold{x}^2\left[\partial_x^j[\partial_jS_k(\bold{x})]\right]=\mu [k]_{q^2}S_k(\bold{x}),
\end{gather*}
which proves the lemma.
\end{proof}

\begin{remark}
Using the same techniques for a general polynomial $P_k\in\cP_k$ yields
\begin{gather*}
\sum_{j,l}\left[\partial_j P_k(\bold{x})\right]^j_lx^l=[k]_{q^2}P_k(\bold{x})+\frac{(q^2-1)}{\mu^2}\bold{x}^2\Delta P_k(\bold{x}).
\end{gather*}
\end{remark}

Equations~\eqref{PizzBessel} and~\eqref{derFK} imply that the quantum sphere integral of the Fourier kernel will yield the f\/irst $q$-Bessel function. This result can be generalized by introducing spherical harmonics in the integration.

\begin{theorem}
\label{FHexp}
For $S_k\in\cS_k$, the following relation holds,
\begin{gather*}
\int_{\mS^{m-1}_q,\bold{x}}S_k(\bold{x})\exp_{\hat{R}}(i\bold{x}|\bold{y}) = 2\mu^{m/2-1} \Big(\Gamma_{q^2}\Big(\frac{1}{2}\Big)\Big)^{m}i^k\frac{J^{(1)}_{\frac{m}{2}+k-1}\big(\frac{1+q}{\mu}\sqrt{\bold{y}^2}|q^2\big)}
{\big(\sqrt{\bold{y}^2}\big)^{\frac{m}{2}+k-1}}S_k(\bold{y}).
\end{gather*}
\end{theorem}

\begin{proof}
First we prove that the relation
\begin{gather}
\label{DSexp}
\Delta_{\bold{x}}^{k+l}S_k(\bold{x})\exp_{\hat{R}}(i\bold{x}|\bold{y})
=i^{k+2l}\frac{[k+l]_{q^2}!\mu^k}{[l]_{q^2}!}\exp_{\hat{R}}(i\bold{x}|\bold{y})\bold{y}^{2l}S_k(\bold{y})+\cdots
\end{gather}
holds, where $\cdots$ stands for terms of the form $P(\bold{x})\exp_{\hat{R}}(i\bold{x}|\bold{y})q^2(\bold{y})$ with $P\in\oplus_{j>0}\cP_j$. In case $k=0$, equation~\eqref{derFK} implies that $\Delta_{\bold{x}}^{l}\exp_{\hat{R}}(i\bold{x}|\bold{y})=(-1)^l\exp_{\hat{R}}(i\bold{x}|\bold{y})\bold{y}^{2l}$ holds. Now we proceed by induction on $k$. Assuming equation \eqref{DSexp} holds for $k-1$ we calculate, using equations~\eqref{Eharm} and~\eqref{Laplx}
\begin{gather*}
\Delta_{\bold{x}}^{k+l}S_k(\bold{x})\exp_{\hat{R}}(i\bold{x}|\bold{y})
= \frac{1}{[k]_{q^2}}\Delta_{\bold{x}}^{k+l}x^j\left[\partial_{x^j}S_k(\bold{x})\right]\exp_{\hat{R}}(i\bold{x}|\bold{y})\\
\phantom{\Delta_{\bold{x}}^{k+l}S_k(\bold{x})\exp_{\hat{R}}(i\bold{x}|\bold{y})}{}
=\frac{1}{[k]_{q^2}}q^{2k+2l}x^j\Delta_{\bold{x}}^{k+l}\left[\partial_{j}S_k(\bold{x})\right]\exp_{\hat{R}}(i\bold{x}|\bold{y})\\
\phantom{\Delta_{\bold{x}}^{k+l}S_k(\bold{x})\exp_{\hat{R}}(i\bold{x}|\bold{y})=}{}
+\frac{[k+l]_{q^2}}{[k]_{q^2}}\mu\Delta_{\bold{x}}^{k+l-1}\left[\partial_{j}S_k(\bold{x})\right]^j_t\partial_{x}^t\exp_{\hat{R}}(i\bold{x}|\bold{y})\\
\phantom{\Delta_{\bold{x}}^{k+l}S_k(\bold{x})\exp_{\hat{R}}(i\bold{x}|\bold{y})}{}
=\cdots+i^{k-1+2l}\frac{[k+l]_{q^2}!}{[l]_{q^2}!}\mu^{k}\exp_{\hat{R}}(i\bold{x}|\bold{y})\bold{y}^{2l}\frac{1}{[k]_{q^2}}\left[\partial_{j}S_k(\bold{y})\right]^j_liy^l.
\end{gather*}
Lemma \ref{Svrexp} then proves that equation~\eqref{DSexp} holds for every $k$. We also used the fact that for $S_p\in\cS_p$, $[S_p]^j_l\in\cS_p$ which follows from the fact that $\partial_x^j$ and $\Delta$ commute. Equation \eqref{DSexp} then yields
\begin{gather*}
\int_{\mS^{m-1}_q,\bold{x}}S_k(\bold{x})\exp_{\hat{R}}(i\bold{x}|\bold{y})\\
 \qquad{} = \sum_{j=0}^\infty\frac{2\left(\Gamma_{q^2}(\frac{1}{2})\right)^{m}}{\mu^{2j+2k}[j+k]_{q^2}!
\Gamma_{q^2}(j+k+\frac{m}{2})}\big(\Delta^{j+k}_{\bold{x}}S_k(\bold{x})\exp_{\hat{R}}(i\bold{x}|\bold{y})\big)(\bold{x}=0)\\
 \qquad{} = i^{k}\frac{2\left(\Gamma_{q^2}(\frac{1}{2})\right)^{m}}{\mu^k}
 \sum_{j=0}^\infty\frac{(-1)^j}{[j]_{q^2}!\Gamma_{q^2}(j+k+\frac{m}{2})}\frac{\bold{y}^{2j}}{\mu^{2j}}S_k(\bold{y}).
\end{gather*}
Comparing this with equation \eqref{Bessel1} proves the theorem.
\end{proof}


This theorem allows to calculate the Fourier transform of an element of $\cP\otimes E_{q^2}(-\beta \frac{\bold{x}^2}{\mu})$ inside one of the irreducible blocks of its $\cU_q(\mathfrak{so}(m))$-decomposition.
\begin{corollary}
\label{BochFour}
The Fourier transform in Definition~{\rm \ref{defFour1}} of a function
\[
S_k(\bold{x})\psi\big(\bold{x}^2\big)\in\cP\otimes E_{q^2}\left(-\beta \frac{\bold{x}^2}{\mu}\right)\subset \overline{\cV}_{[\beta]}
\]
 with $S_k\in\cS_k$ is given by
\begin{gather*}
\overline{\cF}^\pm\big[S_k(\bold{x})\psi\big(\bold{x}^2\big)\big](\bold{y}) = (\pm i)^kS_k(\bold{y})\overline{\cF}^{q,\beta}_{\frac{m}{2}+k-1}[\psi]\big(\bold{y}^2\big)\\
\qquad{}  = (\pm i)^k\frac{1+q}{\mu}\left[\int_{0}^{\sqrt{\frac{\mu}{(1-q^2)\beta}}}d_qr \,r^{m+2k-1}\psi\big(r^2\big)\frac{J^{(1)}_{\frac{m}{2}+k-1}\big(\frac{1+q}{\mu}rt|q^2\big)}{\left(rt\right)^{\frac{m}{2}+k-1}}
\right]_{t^2=\bold{y}^2}S_k(\bold{y})
\end{gather*}
with $r$ and $t$ two real commuting variables.
\end{corollary}

These formulae are the Bochner's relations for the Fourier transform on $\mR^m_q$. For the classical Bochner's relations, see e.g.~\cite{MR1151617}.

This corresponds exactly to the one dimensional case for $\beta=1$. Since $\lim\limits_{m\to1}\mu=1+q$, Corollary~\ref{BochFour} for $m=1$ and $k=0$ reduces to
\begin{gather*}
\frac{\sqrt{1+q}}{2\Gamma_{q^2}(\frac{1}{2})}\int_{-\frac{1}{\sqrt{1-q}}}^{\frac{1}{\sqrt{1-q}}}d_qx\,\psi\big(x^2\big)e_q(ixt) = \left[\int_{0}^{\frac{1}{\sqrt{1-q}}}d_qr \psi(r^2)\frac{J^{(1)}_{-\frac{1}{2}}(rt|q^2)}{(rt)^{-\frac{1}{2}}}\right] \quad\mbox{for}\quad t^2<\frac{1}{1-q}.
\end{gather*}
Corollary \ref{BochFour} for $m=1$, $k=1$ and $S_1=x$ is
\begin{gather*}
\frac{\sqrt{1+q}}{2\Gamma_{q^2}(\frac{1}{2})}\int_{-\frac{1}{\sqrt{1-q}}}^{\frac{1}{\sqrt{1-q}}}d_qx\,x\,\psi\big(x^2\big)e_q(ixt) = i\left[ \! \int_{0}^{\frac{1}{\sqrt{1-q}}} \!\!d_qrr^2 \psi(r^2)\frac{J^{(1)}_{\frac{1}{2}} (rt|q^2 )}{(rt)^{\frac{1}{2}}}\right]t
\quad \mbox{for}\quad t^2<\frac{1}{1-q}.
\end{gather*}
This agrees with
\[
e_q(iu)=\Gamma_{q^2}\Big(\frac{1}{2}\Big)\left(\frac{u}{1+q}\right)^{(1/2)}\left[J_{-\frac{1}{2}}^{(1)}\big(u|q^2\big)
+iJ_{\frac{1}{2}}^{(1)}\big(u|q^2\big)\right],
 \]
 which can be easily calculated.

\section{Properties of the Fourier transform}
\label{behave}



The Fourier transform is determined by its Bochner's relations, see Corollary \ref{BochFour}. The second Fourier transform is immediately def\/ined here by its Bochner's relations.
\begin{definition}
\label{defFour}
The second Fourier transform $\cF^{\pm}_{\mR^m_q}$ on $\mR^m_q$ is a map $\cV_{[\alpha]}\to\overline{\cV}_{[1/\alpha]}$ for each $\alpha\in\mR^+$. For a function $S_k(\bold{y})\psi(\bold{y}^2)\in\cP\otimes E_{q^2}\big({-}\frac{\bold{y}^2}{\mu\alpha}\big)$ the transform is given by
\begin{gather*}
\cF^{\pm}_{\mR^m_q}\big[S_k(\bold{y})\psi(\bold{y}^2)\big](\bold{x}) = (\pm i)^kS_k(\bold{x})\frac{1}{c(\sqrt{\alpha}\gamma)}\cF^{q,\gamma}_{\frac{m}{2}+k-1}\left[\psi\right]\big(\bold{x}^2\big),
\end{gather*}
for an arbitrary $\gamma\in\mR^+$ with
\begin{gather*}
c\big(\sqrt{\alpha}\gamma\big) = \frac{q^{(\frac{m}{2}-1)\frac{m}{2}}}{\Gamma_{q^2}(\frac{m}{2})}\int_0^{\frac{\alpha q^2\gamma^2 }{\mu}\cdot\infty}d_{q^2}u u^{\frac{m}{2}-1}E_{q^2}(-u)=d\left(\gamma\sqrt{\frac{\alpha}{\mu}},\frac{m}{2}-1\right)\\
 \phantom{c\big(\sqrt{\alpha}\gamma\big)}{}
 = d\left(\gamma\sqrt{\frac{q^{-2j}\alpha}{\mu}},\frac{m}{2}+k-1\right)\qquad\mbox{for} \quad j,k\in\mN.
\end{gather*}
\end{definition}
The second Fourier transform does not depend on the choice of $\gamma$ as can be seen from the expressions in the subsequent Theorem~\ref{Fourinv}. From the properties of $c(\sqrt{\alpha}\gamma)$ it is clear that the def\/inition does not depend on which $\cV_\alpha$ the element of $\cV_{[\alpha]}$ is chosen to be in. Although the Fourier transform on each space~$\cV_{[\alpha]}$ is denoted by the same symbol, each Fourier transform should be regarded as an independent operator. In Section~\ref{FourHO} these dif\/ferent transforms will be combined in order to construct the Fourier transform on the Hilbert space corresponding to the harmonic oscillator.


\begin{theorem}
\label{Fourinv}
The Fourier transforms in Definitions~{\rm \ref{defFour1}} and~{\rm \ref{defFour}} are each others inverse, i.e.\ for each $\alpha,\beta\in\mR^+$
\begin{gather*}
\overline{\cF}_{\mR^m_q}^{\mp}\circ\cF^{\pm}_{\mR^m_q} = id_{\cV_{[\alpha]}}\qquad\mbox{and}\qquad
\cF_{\mR^m_q}^{\mp}\circ\overline{\cF}^{\pm}_{\mR^m_q} = id_{\overline{\cV}_{[\beta]}}.
\end{gather*}
\end{theorem}

\begin{proof}
This is a consequence of Corollary~\ref{BochFour} and Theorem~\ref{Hankelinv}. It can also be obtained directly from the relations,
\begin{gather*}
\overline{\cF}^{\pm}_{\mR^m_q}\left[\cL_j^{(\frac{m}{2}+k-1)}\left(\frac{\bold{x}^2}{\alpha\mu}|q^2\right)
S_k(\bold{x})E_{q^2}\left(-\frac{\bold{x}^2}{\alpha\mu}\right)\right](\bold{y})\\
 \qquad{} = (\pm i)^k\alpha^{\frac{m}{2}+k+j}C_j\bold{y}^{2j}S_k(\bold{y})e_{q^2}\left(-\alpha\frac{q^2\bold{y}^2}{\mu}\right),\\
\cF^{\mp}_{\mR^m_q}\left[\alpha^{\frac{m}{2}+k+j}C_j\bold{y}^{2j}S_k(\bold{y})e_{q^2}\left(-\alpha\frac{q^2\bold{y}^2}{\mu}\right)\right](\bold{x})\\
 \qquad{} = (\mp i)^k\cL_j^{(\frac{m}{2}+k-1)}\left(\frac{\bold{x}^2}{\alpha\mu}|q^2\right)S_k(\bold{x})E_{q^2}\left(-\frac{\bold{x}^2}{\alpha\mu}\right),
\end{gather*}
with $C_j=\frac{q^{2(j+1)(j+\frac{m}{2}+k)}}{[j]_{q^2}!\mu^j}$. These follow immediately from the calculations before Theorem~\ref{Hankelinv}.
\end{proof}
